\documentclass[12pt]{article}

\usepackage{amsmath,amssymb,amsthm}
\usepackage{hyperref}

\newtheorem{theorem}{Theorem}
\newtheorem{lemma}{Lemma}
\newtheorem{definition}{Definition}
\newtheorem{remark}{Remark}

\title{A Note on the Asymptotic Distribution of the Weighted Median Estimator}
\author{Test Author}
\date{}

\usepackage{color}
\newcommand{\grammarcheck}[2]{\textcolor{red}{#1}\,\textcolor{red}{[#2]}}
\newcommand{\techcheck}[1]{\textcolor{blue}{[#1]}}

\begin{document}
\maketitle

\section{Introduction}

The weighted median is a robust estimator widely used in high-dimensional regression and quantile estimation. While its consistency has been studied extensively, the asymptotic normality remains not fully \grammarcheck{understoond}{typo: ``understood''}. In this paper, \grammarcheck{we derives}{subject-verb agreement: ``we derive''} the limiting distribution of the weighted median under mild regularity conditions and illustrate its applications to heteroskedastic regression.

Compared to the sample mean, the weighted median is more resilient to outliers and heavy-tailed noise, \grammarcheck{makes it particular useful}{awkward: consider ``making it particularly useful''. Dangling participle --- attach to main sentence or start a new one.} in financial and biomedical applications where data quality \grammarcheck{can not}{style: ``cannot'' is standard in academic writing} be guaranteed.

\section{Notation and Preliminaries}

Let $X_1, \ldots, X_n$ be i.i.d. random variables from a distribution $F$ with density $f$. Let $w_1, \ldots, w_n$ be non-negative weights satisfying
\begin{equation}
    \sum_{i=1}^n w_i = 1.
\end{equation}
The weighted median $\hat{\mu}_n$ is defined as
\begin{equation}
    \hat{\mu}_n = \inf\left\{ x : \sum_{i=1}^n w_i \mathbf{1}\{X_i \leq x\} \geq \frac{1}{2} \right\}.
\end{equation}

We denote the true median as $\mu_0 = F^{-1}(1/2)$. The following standard conditions are needed for the main theorem.

\begin{definition}
The distribution $F$ is said to be \textit{median-regular} if $F$ is continuously differentiable at $\mu_0$ and $F'(\mu_0) = f(\mu_0) > 0$.
\end{definition}

\section{Main Results}

\begin{theorem}\label{thm:main}
Let $X_1, \ldots, X_n$ be i.i.d. with a median-regular distribution $F$, and let $w_i = w(X_i)$ where $w(\cdot)$ is a bounded positive function with $\mathbb{E}[w(X)] = 1$. Then, as $n \to \infty$,
\begin{equation}
    \sqrt{n}(\hat{\mu}_n - \mu_0) \xrightarrow{d} N\left(0, \sigma^2\right),
\end{equation}
where
\begin{equation}
    \sigma^2 = \frac{\mathbb{E}[w(X)^2 \mathbf{1}\{X \leq \mu_0\}]}{4 f(\mu_0)^2 \mathbb{E}[w(X)]}.
\end{equation}
\end{theorem}

\begin{proof}
Define the empirical process
\begin{equation}
    G_n(x) = \sum_{i=1}^n w_i \mathbf{1}\{X_i \leq x\}.
\end{equation}
By the law of large numbers, for each fixed $x$,
\begin{equation}
    G_n(x) \xrightarrow{p} \mathbb{E}[w(X) \mathbf{1}\{X \leq x\}] \equiv G(x).
\end{equation}
Note that $G(\mu_0) = 1/2$ by the normalization condition on $\mu_0$. The function $G$ is differentiable at $\mu_0$ with $G'(\mu_0) = w(\mu_0) f(\mu_0)$.
\techcheck{Under the deterministic weight function $w(\cdot)$, this step is correct since $\frac{d}{dx}E[w(X)1\{X\leq x\}] = E[w(X)|X=x]f(x) = w(x)f(x)$. However, note the inconsistency on line~78 where $G'(\mu_0)$ is later claimed to equal $f(\mu_0)$.}

The estimator $\hat{\mu}_n$ solves $G_n(\hat{\mu}_n) = 1/2$. By the Bahadur-Kiefer representation,
\begin{equation}
    \sqrt{n}(\hat{\mu}_n - \mu_0) = \frac{\sqrt{n}(G_n(\mu_0) - G(\mu_0))}{G'(\mu_0)} + o_p(1).
\end{equation}
\techcheck{The Bahadur-Kiefer representation with $o_p(1)$ remainder typically requires $f$ to be continuous in a \emph{neighborhood} of $\mu_0$ (not just differentiable at a point), and often needs a second derivative or Lipschitz condition on $G$. The stated ``continuously differentiable at $\mu_0$'' condition may be insufficient to guarantee the $o_p(1)$ bound.}

The central limit theorem gives
\begin{equation}
    \sqrt{n}(G_n(\mu_0) - G(\mu_0)) \xrightarrow{d} N(0, \tau^2),
\end{equation}
where $\tau^2 = \operatorname{Var}(w(X) \mathbf{1}\{X \leq \mu_0\})$. The conclusion follows by Slutsky's theorem, and the variance simplifies to the stated expression after noting that $G'(\mu_0) = f(\mu_0)$ since $\mathbb{E}[w(X)] = 1$.
\techcheck{Critical inconsistency: $G'(\mu_0)$ was derived as $w(\mu_0)f(\mu_0)$ on line~67, but here it is claimed to equal $f(\mu_0)$. The condition $\mathbb{E}[w(X)]=1$ does \emph{not} imply $w(\mu_0)=1$, so this simplification is unjustified. The variance expression $\sigma^2$ in Theorem~1 may need to retain the factor $w(\mu_0)$ or be re-derived.}
\end{proof}

\begin{remark}
The assumption $\mathbb{E}[w(X)] = 1$ is made without loss of generality, since the weights can always be renormalized. However, this normalization must be done carefully when weights depend on estimated parameters.
\end{remark}

\section{Simulations}

We conduct a small simulation study to verify the asymptotic approximation \grammarcheck{in finite sample}{missing plural: ``in finite samples''}. Data is generated from a standard Cauchy distribution, and weights are set proportional to $|X_i|$. Figure~\ref{fig:qq} \grammarcheck{display}{subject-verb agreement: ``displays''} the QQ-plot of the standardized estimator against the standard normal.
\techcheck{Simulation setup: standard Cauchy data with weights $\propto |X_i|$. For Cauchy, $\mathbb{E}|X| = \infty$, so the weights are unbounded almost surely --- this violates the bounded-weight assumption of Theorem~1. The deviation from normality at $n=100$ may reflect assumption violation rather than genuine finite-sample behavior. Consider using a distribution with finite moments (e.g.\ $t_3$) or bounded weight functions.}

The simulation \grammarcheck{results confirms}{subject-verb agreement: ``results confirm''} the theoretical prediction for $n = 1000$, but reveal \grammarcheck{noticeable deviation}{plural: ``noticeable deviations'' (or ``a noticeable deviation'')} from normality when $n = 100$, particularly in the tails. This is expected because the weighted median has heavier tails than the normal distribution for small sample sizes.

\section{Discussion}

In this paper, the asymptotic normality of the weighted median estimator under data-dependent weights was established. A key limitation is that our \grammarcheck{theorem require}{subject-verb agreement: ``theorem requires''} the weights to be bounded, which excludes certain important cases such as inverse-variance weights. Extension to unbounded weight functions and to dependent data are \grammarcheck{left as future work}{preposition: ``left for future work''}. The simulation study \grammarcheck{also suggest}{subject-verb agreement: ``also suggests''} that the normal approximation may be poor unless the sample size is quite large, \grammarcheck{especialy}{typo: ``especially''} for heavy-tailed distributions.

\end{document}
